\clearpage\thispagestyle{plain}
\phantomsection\addcontentsline{toc}{chapter}{致谢}
\markboth{致谢}{致谢}
\vspace*{8mm}
\begin{center}{\zihao{2}\bfseries\color{FrontBlue}致\quad 谢}\end{center}
\vspace{10mm}
{\large\bfseries\color{FrontBlue}联系方式}\par\smallskip
这本复习册记录了对线性代数概念、计算方法和几何结构的整理与思考。学习中若发现定义、推导、题目解答或排版上的问题，欢迎附上章节、例题编号和具体疑问交流。每一条清楚的反馈，都能帮助这本书变得更准确、更易读。
\begin{itemize}[leftmargin=1.7em,itemsep=5pt]
\item 知乎：\href{https://www.zhihu.com/people/black-despair}{\textcolor{FrontBlue}{\nolinkurl{https://www.zhihu.com/people/black-despair}}}
\item QQ邮箱：\href{mailto:3547908579@qq.com}{\textcolor{FrontBlue}{\nolinkurl{3547908579@qq.com}}}
\end{itemize}
\vspace{8mm}
{\large\bfseries\color{FrontBlue}参考资料与感谢}\par\smallskip
感谢教材、课程讲义与专题材料的编写者。严谨的定义、清晰的证明和各有侧重的解题方法，为本册梳理知识、比较多解和补充几何解释提供了启发。
\begin{enumerate}[leftmargin=1.8em,itemsep=6pt]
\item Gilbert Strang，\textit{Introduction to Linear Algebra}，第五版（中文译本）：向量、子空间、正交投影与矩阵分解的几何观点。
\item 杨启哲，《线性代数》课程讲义：线性空间、线性映射及核像结构的系统组织。
\item 冯建轩，《线性代数大讲义》：考研知识脉络与常见题型的归纳。
\item 何军华，《逆矩阵》专题课件：逆矩阵的定义、等价刻画、初等变换与矩阵方程。
\item 《\(n\)阶行列式的计算方法》：展开、递推、拆项、加边与结构计算等专题方法。
\end{enumerate}
同时感谢认真演算、追问条件并愿意指出问题的读者。一个被发现的反例、一处补齐的前提、一份能够复算的解答，都会使知识之间的联系更清楚。愿这本复习册能陪伴你把计算练熟，也把概念想明白。
\vspace{7mm}
\begin{flushright}雪落长安无声\par 2026.10.31\end{flushright}
\vfill
{\footnotesize\color{black!55}封面书页照片：freestocks.org，\href{https://commons.wikimedia.org/wiki/File:Open_book_(26421751605).jpg}{\textcolor{black!55}{Open book}}，CC0；封面寄语为本册自拟。}
